The market at the Slate Row end went thin in layers, the way a kiln goes cold first at the door. The people who had business in the narrow street drifted toward it and left their smells behind them — barley, wet straw, the sour woodsmoke of a barrow standing with its shaft up — and what they left was not emptiness but the shape of what had just gone, which is a thing a man learns to read if he has ever watched a ward empty out after a shift and counted the beds that would not be filled again that night. \person{Sella Vail} set the basket down beside a post with a notice nail in it, not where the slate lay open, and the post gave a finger's width of shade at the foot, the only shade at that end of \place{the Cross}. \person{Corin} lay with his cheek to the cold iron of the post and felt the market come through it in small tremors, the same tremors a kiln wall gives when the wheel takes a heavy piece, and counted the tremors because there was nothing else at that end of the square to count.

A table stood at the fringe where the Court's business was done in the open air before it was carried into \place{Slate Row}, which the market spoke of as a direction and not a place that anybody at this end meant to go that day. A clerk sat behind it. He was a narrow man with a cuff gone soft at the wrist from much washing and ink in the creases of two fingers, and the skin of those fingers was the colour of other men's fingers, which Dan noticed because he noticed hands before he noticed faces and was always relieved, in a country full of officers, when a hand looked ordinary. He had a pile of slates to one side and a rag for the chalk, and he worked with the unhurried care of a man who has entered a great many names and has never once been paid to hurry.

When a woman came to the table with a question in her face, the clerk put his hand on the top slate before he looked at her. It was not rudeness. It was the order in which the table had taught him to do things, the same order \person{Alda Reed} had kept in a room three doors up the yard, where her hand went to the birth-slate on the bench before it went to the child. Dan knew the order of reaching before he knew what was being reached for; he had learned it in four days the way other men learn a trade, by watching hands that did not know they were being watched. The woman's question waited while the slate was set straight. Then the clerk looked at her and said, in a soft voice without any edge in it, “It's entered as it stands. There's nothing in it the household need correct.”

He did not say I. He did not say the district, the way the man on the block had said it; he said it and the household, and the it stood for a line of writing that had been put down somewhere before either of them had opened a mouth. He was an officer and he had a reason, and the reason was in the way he laid the rag down: a man who has to get a morning's names into slates before the market wears out, and has found the quickest way is to touch the slate first and explain afterwards. Dan had made the same bargain once in a hospital in \place{Chicago}, in a room with a long table where the chart went under the hand before the patient's face did, and had never thought of it as a bargain until he was on the other side of one. He set the thought with the others and did not make anything of it. The thoughts that came in threes were the ones to keep; the one that tried to be a fourth was usually his own.

The upper line on the slate was in a strange neat hand.

That was all he could have said about it in \lore{Grathic}, and he could not have said it in Grathic, so he said nothing and read the shape instead. The marks were small and upright and very even, with the corners squared off the way a potter squares a lip before the fire, and there was a space between the first cluster and the second where the writing had stopped and started again. A stop, not a fault. The hand that made it had lifted and set down again in the same place each time, as if the writer had been counting and had put the count down between the strokes. Dan had seen that stop once at the edge of a flare light in the country he had come from, and once on a slate lying flat on a bench under a woman's hand in a room where a child had just been born. Now he saw it a third time, in a market, on a line that had been set down before anybody there had thought to look at the face it belonged to.

He kept the three side by side the way he kept his numbers, and did not make a fourth. He could not read a word of the line. He could read a hand, which is a different thing and, in his trade, the older one: the pressure, the lift, the small deciding pause before a stroke that a man does not know he makes. The pause was the whole of what he had, and it was no more use to him here than a pulse in a throat two paces off would have been. Nobody at the table looked into the basket. Nobody in the square remarked on the line, or on the fact that it sat above the others as if it had been there before the day began, which was the way the district had of putting a thing where everyone could see it and nobody would think to ask about it. The cold got into his mouth under the wool; a thread of drool went from his lip to the fold and stayed. He tried to push the fold back with the hand that had once reduced a shoulder in the dark, and the hand was a small soft thing outside the basket with milky nails that closed on air an inch from his cheek and came away. He had put a thousand shoulders back. He could not keep a fold of wool on his own chin. The head went where it wanted, and past a certain point it took a small wet sound with it, and the sound got the wool put right by a hand that did not look and was gone again before he could be ashamed of it.

\scenebreak

\person{Berr Clay} came down the fringe from the direction of the moot, wiping his thumb on his apron before he reached the table, though he was not going to touch the slate at the table; he had touched his own slate that morning and had set it straight twice already, and the wipe was for the hand and not for the stone, which Dan understood better than most men in the square. The red wool cord of his office lay folded across his chest, and he took it off only when he went home, folding it into a drawer, so that a man could see the office coming before he saw the man. He smelled of hot clay and lead, and when he leaned down toward the basket to speak over Sella's shoulder the lead came on his breath, sweet and a little sick, the taste of a man who had been tasting his own glaze since morning and had forgotten he was doing it. The care he spent on a clean slate was the care of a man who keeps a thing he could produce if he ever chose to; Dan had known men who kept such things, and knew the certainty in the hand before he knew what the hand held.

“Sella,” he said. “Your widow's mark is confirmed by the factor at \place{Greybank}. That is a kindness, and kindness is entered.”

He said it the way he had said the figure at the block: low, and to the woman's shoulder rather than to her face, with the grammar of the district inside it so that no part of the sentence held a man who could be answered. The factor at Greybank had confirmed it; the confirmation was a kindness; the kindness was entered. There was nobody in the sentence who had decided to be kind, and so nobody who could be thanked, and so no thanks that a woman could refuse. Dan had heard the shape of it before, in a room with good light and a window that would not open, from a man with a file under his arm who had said a chart was a chart and that nothing had happened in the room that the chart did not already know about. \memory{That's what a chart is for.} The sentence had come warm and reasonable and had closed a matter that had never been open, and Dan had carried it out of that room and into the rest of his life without once saying what it had done. Here it had gone outward and was standing in a market in a man's mouth with lead on his breath. It was the first time he had heard the grammar outside his own head. He had kept it inside for two days. Now it had a voice and a green ring at the corner of a mouth, and it was no less reasonable for being in the open, which was the frightening part and the part he would have recognised anywhere, in either country.

Sella did not answer him about the mark. She stood with her hand on the post, and the little-door burn on the back of her left hand — two straight sides and a lintel and a threshold — closed on the iron and opened again, and she put her hand to the fold of wool at his chin instead, pulling it up over his ear without looking down at him.

“Keep it off his ears,” she said to the wool, or to the post, or to the cold coming off the river. “It's the river wind.”

That was the whole of her answer, and it was an instruction to a thing that could not answer back, given in the form she had used since she took him in: the house would keep him, the house did not explain, and the boy's ears were to be covered because the wind was off the river. Berr stood a moment with his mouth open on the next part of what he had come to say, and Dan watched the not-saying happen the way he had once watched a resident keep his eyes on his shoes at a review of a chart. There was a space in the man's sentence, and the space had a shape. It was the shape of a debt left where it lay, against a firing of tiles that had gone into the floor of this market years ago and had never been entered on the slate Berr kept so clean. His care not to call it was another kindness entered nowhere, for a debt uncalled can never be paid and never be settled and never be done with; it sits where it is as solid as the hollow seat in the house up the yard that Sella will not let anybody fill. Dan knew the shape of an uncalled thing. He had lived beside one for three days in a house where the ash in the corner was the thickest in \place{Low Kiln}, and he had left one of his own across an ocean in a coat with no phone in it, and set side by side they had the same shape, a shape with a gap where the fifth thing should have been.

Berr did not call the debt. He set his thumb to his apron again, one short dry stroke, and his eyes went to the cord at Sella's throat where the widow's mark lay under her collar — a thumb-length of lead-glazed clay with the Vail kiln scratched into it, which she touched before she opened the kiln and at no other time — and he did not touch it and did not ask to see it. Then he straightened. “The figure stands as read,” he said. “There'll be a slate at the watch if the house wants it again.”

He had come to confirm a kindness and he had confirmed it, and he went back toward the moot with the lead going off him and the red cord folded square, leaving the space where the debt was as carefully unentered as he had found it. Dan lay with his ears uncovered where Sella's hand had put the wool and watched the man go, and knew the shape of what had passed between them, and had no word in Grathic for any part of it and nobody to say it to if he had.

\scenebreak

Alda Reed came to the edge of the market with a bundle of birth-slates under one arm and her grey wool shawl pulled tight across her shoulders, which was the way she wore it in every season. The little tin on its cord was inside her bodice, where she touched it when a question was asked and never opened it in company. She had been keeping the birth-slates since the market opened, setting them in a row on a bench at the fringe and turning the writing down when a mother came near, the way she turns a slate flat on the bench in a birth room before she reaches for the child. The slates were flat grey things with chalk in their creases, and she handled them with the steady hands of a woman who has lifted four hundred children and never dropped one or a slate.

She stopped beside the basket to set the bottom slate straight, and her hand went once to the tin inside her bodice, a small pressed movement that a man would not notice unless he had been watching for it, which Dan had been watching for it since the day he was born. He had seen her touch that tin in a birth room and again over a bundle going up the yard, and here was a third time, and it was not yet a number that meant anything, only a thing that happened when a question came near. There was no question now. The wool had ridden back off his wrist where Sella's hand had pulled it at his ear, and the edge of the birth-flower showed there, the red branching end of the mark that ran down his forearm and stopped short of the joint — a flower's edge the colour of a line drawn under a word.

Alda looked at it. Dan saw her look, and saw what her hand did before she knew she had done it: two fingers crossed at the corner of the mark, the sign her mother had made and she had made over four hundred children, and did not know she was making it in front of a market full of people. It was done and gone in the time it takes to settle a shawl, and then she covered it the way she covered everything, with a date.

“Stone seven,” she said, to nobody, to the slate, to the row of flat grey things under her arm. “The account went up on the eighth, before the paper went out. A child born on the seventh is entered on the eighth, and that's the whole of it.”

The date arrived too fast, the way her dates always did, and it arrived with no question in the air to answer. Dan lay under it and did not move. Nobody at the bench commented on the sign or the date or the red edge at his wrist. A woman bargaining over short weight called a figure across the market and another woman called a different figure, and the two went on standing side by side without either troubling to make them agree. \lore{The Hearth Drawing} was two days off, and \place{Low Kiln} said \lore{the Five} were three days off, and the market went on arranging itself around the things it did not read, which was most of what was in it.

Sella put her hand back to the fold at his wrist. The little-door burn closed on the wool — lintel, threshold, two straight sides — and the fold went up over the red edge, and the hand was gone again as if it had never been there, which was the way she did everything, leaving no trace of the feeling and only the practical fact of a covered wrist. A pace away, at the Court's table, the clerk set his hand on the slate and drew the rag across it and turned it straight for the next name. Her hand and his hand were both careful hands, both closing on the same child in two different accounts: the house's private roll entered in a careful hand that would not be shown to a man come down from Slate Row, and the Court's slate with its neat upper line set down before anybody had seen the face it belonged to. Neither account had been shown to him. Neither had asked whether he wanted to be entered in it. Dan lay with his cheek to the cold iron of the post and his own small hand outside the fold with milky nails, and tried to lift it toward the wool; the hand went up an inch and came down, and the two careful hands went on entering him on either side of a space he could not fill.
